% Our system draws on four lines of prior work: maritime semantic segmentation, monocular depth estimation, multi-object tracking, and collision-risk modeling.

\subsection{Maritime Semantic Segmentation}

Fully-convolutional architectures such as U-Net~\cite{ronneberger2015unetconvolutionalnetworksbiomedical} and the DeepLab family~\cite{chen2017rethinkingatrousconvolutionsemantic} established atrous convolution and spatial-pyramid pooling as standard tools for enlarging the effective receptive field, and transformer-based SegFormer~\cite{xie2021segformersimpleefficientdesign} later matched or exceeded these CNN baselines with a lighter hierarchical-attention design. In the maritime domain, WaSR~\cite{bovcon2021wasr} and MODD~\cite{Kristan2015} adapted this line of work to cope with wakes, reflections, and sun glitter, and the LaRS~\cite{zust2023larsdiversepanopticmaritime} and WaterScenes~\cite{yao2024waterscenesmultitask4dradarcamera} datasets have since consolidated training and evaluation data for ASV perception, with LaRS serving as our primary source. Nearly all of this work, however, optimizes pixel-level mIoU as the end goal, a metric that rewards precise boundaries but is silent on whether a hazard was detected at all. We instead treat segmentation as one stage of a downstream warning pipeline: evaluation is re-scaled around obstacle recall and inference latency, and SegFormer is extended with a boundary-prediction head, absent from the original design, to recover instance separation from a semantic-only backbone.

\subsection{Monocular Depth Estimation}

Monocular depth estimation has moved from dataset-specific regressors toward generalist foundation models. MiDaS~\cite{ranftl2020robustmonoculardepthestimation} showed that mixing depth datasets during training enables zero-shot relative-depth transfer, Depth Anything~\cite{yang2024depthanything} scaled this idea to tens of millions of unlabeled training images, and Depth Anything V2~\cite{yang2024depthv2} sharpened object boundaries with synthetic data and a larger teacher network. Diffusion-based methods such as Marigold~\cite{ke2025marigoldaffordableadaptationdiffusionbased} and DepthFM~\cite{gui2024depthfmfastmonoculardepth} trade inference speed for finer detail recovery. All of these models are trained primarily on road-driving and indoor imagery~\cite{gaidon2016virtualworldsproxymultiobject}; applied directly to open water, they show a consistent domain gap, with the horizon saturating and metric scale drifting away from physical units. Rather than retraining on scarce maritime depth data, we keep the pre-trained Depth Anything V2 representation intact and correct this drift with a single global scaling factor fit against radar ground truth, restoring near-range metric accuracy in the close-up range that matters most for collision warning.

\subsection{Multi-Object Tracking}

SORT~\cite{bewley2016simple} is a widely-used online tracker that combines a constant-velocity Kalman filter on 2D bounding boxes with Hungarian matching on IoU. DeepSORT~\cite{wojke2017simpleonlinerealtimetracking} augments SORT with a learned appearance embedding to reduce identity switches through occlusion. We adopt the SORT formulation for its low latency and minimal training requirements, but extend its 7-dimensional Kalman state to nine dimensions so that metric depth and its rate of change are tracked jointly with bounding-box dynamics.

\subsection{Collision-Risk Modeling}

Two lines of thought have shaped this area. The first is formal safety, exemplified by Mobileye's Responsibility-Sensitive Safety (RSS) framework~\cite{shalevshwartz2018formalmodelsafescalable}, which derives longitudinal and lateral safe distances from reaction time and braking capability, and by the ISO 22839 standard for forward collision mitigation~\cite{iso22839}. The second relies on domain-specific heuristics: the ship-domain concept of Fujii and Tanaka~\cite{fujii1971traffic} sizes a safe region around a vessel from its beam and operating speed, and recent ASV work adapts COLREGs rules to reactive collision avoidance~\cite{qu2022real}. Both typically collapse risk to a binary safe/unsafe decision and assume range measurements from active sensors. Our model keeps RSS's geometric decomposition into a forward-range term and a lateral-excess term, but replaces the kinematic safe-distance equations with a continuous exponential-decay score, weighted by obstacle class and modulated by closing velocity, so that approaching and receding obstacles at the same range are no longer treated as equally risky.
